\mode*
\section{Introduction}

Version control is a standard part of computing curricula: scholars agree
that students must be prepared for distributed, collaborative workflows
\autocite{IsomottonenCochez2014}, and Git is the system they meet.
Yet Git is famously confusing for learners.
The confusions are specific and documented: students conceptualise a branch
as a directory and try to move into it with \texttt{cd}
\autocite{IsomottonenCochez2014}; multi-institution data from the GitKit
learning environment catalogues misconceptions in basic GitHub workflows
\autocite{Postner2026GitMisconceptions}.
Notably, the difficulty is not only the learner's: an analysis of Git's
design explains well-known difficulties as \emph{concept
misfits}---underlying concepts failing to meet their intended purposes
\autocite{DeRossoJackson2016Gitless}.
When the tool itself is conceptually crooked, teaching must make the
underlying model discernible rather than leave students to reverse-engineer
it from command behaviour.

\mode<presentation>{%
\begin{frame}
  \begin{block}{Git is confusing---specifically}
    \begin{itemize}
      \item Branch conceptualised as a directory: students try to
        \texttt{cd} into it \autocite{IsomottonenCochez2014}.
      \item Misconceptions in basic workflows catalogued across
        institutions \autocite{Postner2026GitMisconceptions}.
      \item Partly the tool's fault: concept misfits in Git's design
        \autocite{DeRossoJackson2016Gitless}.
    \end{itemize}
  \end{block}
\end{frame}%
}

In this work we apply the approach of our companion paper on introductory
programming \autocite{VTProgMisconceptions}: we treat documented
misconceptions as the phenomenographic observations for a variation-theory
analysis \autocite{NCOL}, identify the critical aspects learners must
discern, and outline patterns of variation to teach them.
We cover Git's commit-graph model, the staging area, branches, and remotes.

We ask the following research questions.

\begin{frame}
\begin{restatable}{question}{rqmisconceptions}\label{rq:misconceptions}
  Which misconceptions and difficulties with Git---its commit model, the
  staging area, branches, and remotes---does the literature document?
\end{restatable}

\begin{restatable}{question}{rqaspects}\label{rq:aspects}
  Which critical aspects of version control with Git do these
  misconceptions and difficulties point to?
\end{restatable}

\begin{restatable}{question}{rqpatterns}\label{rq:patterns}
  Which patterns of variation can make these critical aspects discernible
  to learners?
\end{restatable}
\end{frame}

Our contributions are:
a review of the literature on learners' difficulties with Git, answering
\cref{rq:misconceptions};
a variation-theory analysis identifying the critical aspects, answering
\cref{rq:aspects};
and a set of patterns of variation for teaching Git, answering
\cref{rq:patterns}.
